%%% FORMATO E IDIOMA %%%
    \documentclass[letterpaper,DIV=15,12pt]{scrartcl}
    \usepackage[spanish,mexico,shorthands=off,es-noindentfirst]{babel}
    \usepackage{scrlayer-scrpage}
        \clearpairofpagestyles
        \ihead{\footnotesize \textit{Teoría de Conjuntos III}}
        \ohead{\footnotesize \textit{2026-2}}
        \cfoot{\normalfont\thepage}
        \addtokomafont{title}{\bfseries \rmfamily}
        \setlength{\headsep}{5pt}
        \newpairofpagestyles{beginstyle}{
            \clearpairofpagestyles
            \KOMAoptions{headsepline=false}
            \cfoot{\footnotesize \pagemark}
            \cfoot{\normalfont\thepage}
        }
        \addtokomafont{section}{\raggedleft\rmfamily}
        \addtokomafont{subsection}{\raggedleft\rmfamily}
        \addtokomafont{subsubsection}{\raggedleft\rmfamily}
        \setlength{\headsep}{8pt}
        \setlength{\footskip}{20pt}
        \setlength{\textheight}{600pt}
%%% PAQUETES %%%
    \usepackage[style=mexican]{csquotes}
    \usepackage{ifthen}
    \usepackage[shortlabels]{enumitem}
        \setenumerate[1]{label=\MakeLowercase{\roman*}), ref=\roman*}
        \setenumerate[2]{label=\MakeLowercase{\alph*}), ref=\alph*}
    \usepackage{mathtools}
    \usepackage[predefined={parent=section}]{keytheorems}
    \usepackage[T1]{fontenc}
    \usepackage{stix2}
    \usepackage[hidelinks]{hyperref}
	\renewcommand{\qedsymbol}{$\blacksquare$}
	\addto\captionsspanish{\renewcommand{\proofname}{\normalfont\bfseries Demostración}}	
	\newenvironment{solucion}{\renewcommand{\qedsymbol}{$\boxtimes$}\begin{proof}[\normalfont\bfseries Solución]}{\end{proof}}
           
%%% C O M A N D O S %%%
    \renewcommand{\emptyset}{\varnothing}
    \NewDocumentCommand{\mb}{O{}}{\cramped{V^{(B)}_{#1}}}      %V^(B)
    \NewDocumentCommand{\mg}{O{}}{\cramped{V^{(\Gamma)}_{#1}}} %V^(Γ)
    \DeclarePairedDelimiter{\vb}{\lBrack}{\rBrack}             %|| - ||
    \DeclareMathOperator{\dom}{dom}                            %dominio
    \DeclareMathOperator{\cod}{cod}                            %codominio
    \DeclareMathOperator{\im}{im}                              %imagen
    \newcommand{\id}{\mathrm{id}}                              %identidad
    \DeclareMathOperator{\stab}{stab}                          %estabilizador 
    \DeclareMathOperator{\fn}{Fn}                              %Fn(-,-) 
    \DeclareMathOperator{\arriba}{\uparrow}                    %arriba
    \DeclareMathOperator{\abajo}{\downarrow}                   %abajo 
    \DeclareMathOperator{\inte}{int}                           %interior
    \DeclareMathOperator{\cer}{cl}                             %cerradura
    \NewDocumentCommand{\est}{m}{\mathfrak{#1}}                %estructura  
%conjuntos
\newcommand{\st}{}                                              
\newcommand{\SetSymbol}[1][]{%
\nonscript\:#1\vert
\allowbreak
\nonscript\:
\mathopen{}}
\DeclarePairedDelimiterX\set[1]\{\}{%
\renewcommand\st{\SetSymbol[\delimsize]}
#1
}
%% D O C U M E N T O %%

\begin{document}

\thispagestyle{plain}

    \begin{center}
        {\fontsize{30}{60}\rmfamily \textbf{Conjuntos simétricos}} \\ \vspace{.2cm} Teoría de Conjuntos III, 2026-2
    \end{center}
    \begin{flushright}
        \footnotesize \hfill Profesor: Luis Jesús Turcio Cuevas.\\
        \hfill Ayudante: Hugo Víctor García Martínez.
    \end{flushright}

\section{Axioma de elección}

Como ya hemos visto, en el método de forcing se construyen modelos de
\textsc{zfc}, por lo que dar un modelo de la negación del axioma de elección
usando este método es imposible. Sin embargo, es posible hacer una modificación
que permita construir modelos de \textsc{zf} donde no se satisfaga el axioma de
elección. 

Lo que haremos es considerar el modelo \(\mb\) y quedarnos con sus elementos
\textit{simétricos}. Así, lo primero es definir qué es ser simétrico. Para esto
empezamos recordando algunos conceptos de teoría de grupos.

\begin{definition}
Sean \(G\) un grupo y \(X\) un conjunto. Una \emph{acción} de \(G\) sobre \(X\)
es una función \(\alpha\colon G\times X\to X\) que satisface:
\begin{enumerate}
    \item \(\alpha(e,x)=x\) para toda \(x\in X\), donde \(e\) es el elemento
    neutro de \(G\). 
    \item \(\alpha(g,\alpha(h,x))=\alpha(gh,x)\) para toda \(g,h\in G\) y 
    \(x\in X\).
\end{enumerate}
\end{definition}

Usaremos la notación usual de \(g\cdot x\) para denotar \(\alpha(g,x)\). Con
esta notación, las condiciones anteriores se pueden escribir como \(e\cdot x=x\)
y \(g\cdot(h\cdot x)=(gh)\cdot x\).

Dada una acción \(G\times X\to X\) es claro que si \(g\in G\) entonces la
función \(x\mapsto g\cdot x\) es una biyección de \(X\) en sí mismo. En nuestro
caso queremos tomar una acción de un grupo \(G\) en un álgebra de Boole completa
\(B\). En esta situación necesitamos que la acción respete la estructura de
\(B\), es decir, que \(g\cdot-\colon B\to B\) sea un morfismo de álgebras de
Boole completas. Cuando suceda lo anterior diremos que \(G\) actúa por
automorfismos de \(B\).

El siguiente paso es extender una acción por automorfismos de \(G\) sobre un
álgebra de Boole completa \(B\) a una acción sobre una estructura \(B\)-valuada.

\begin{definition}[label=def:aest]
    Sean \(B\) un álgebra de Boole completa y \(\est{A}\) una estructura
    \(B\)-valuada. Una acción de un grupo \(G\) sobre \(\est{A}\) es una acción
    de \(G\) sobre \(B\) por automorfismos y una acción sobre \(A\) tales que
    para cualquier fórmula atómica \(\varphi\) se satisface
    \begin{equation*}
        \vb{\varphi(ga)}=g\cdot\vb{\varphi(a)}.
    \end{equation*}
\end{definition}

Para diferenciar las acciones sobre \(B\) y sobre \(A\), en la definición
anterior, usamos \(g\cdot b\) para denotar la acción sobre \(B\) y \(ga\) para
denotar la acción sobre \(A\).

\begin{proposition}[label=prop:afor]
    Sea \(G\times\est{A}\to\est{A}\) una acción.
    Si \(\varphi\) es una fórmula, \(a\in A\) y \(g\in G\), entonces
    \begin{equation*}
        \vb{\varphi(ga)}=g\cdot\vb{\varphi(a)}.
    \end{equation*}
\end{proposition}
\begin{proof}
    La demostración es por inducción sobre la formación de fórmulas.
    Las atómicas se tienen por definición de acción. Todos los conectivos y
    cuantificadores se hacen de la misma forma, así que sólo haremos el caso del
    existencial.
    \begin{align*}
        \vb{\exists x\varphi(gx)} &= \bigvee_{a\in A}\vb{\varphi(ga)}\\
        &= \bigvee_{a\in A}g\cdot\vb{\varphi(a)} && \text{hipótesis de inducción}\\
        &=g\cdot\bigvee_{a\in A}\vb{\varphi(a)} && \text{acción por automorfismos}\\
        &=g\cdot\vb{\exists x\varphi(x)}. \tag*{\qedhere}
    \end{align*}
\end{proof}

Ahora, dada una acción por automorfismos \(G\times B\to B\) definimos una
función \(G\times\mb\to\mb\) por recursión sobre la relación \enquote{estar en
el dominio de} como sigue
\begin{equation*}
    \dom(gx)=\set{gt\st t\in\dom(x)}, \quad gx(gt)=g\cdot x(t).
\end{equation*}

\begin{proposition}
    La función anterior define una acción sobre \(\mb\).
\end{proposition}
\begin{proof}
    Todas las demostraciones se harán por inducción.
    Primero veamos que actuar con el neutro del grupo es la identidad. Para esto
    notamos
    \begin{alignat*}{2}
        \dom(ex)&=\set{et\st t\in\dom(x)} &\qquad ex(et)&=e\cdot x(t)\\
        &=\set{t\st t\in\dom(x)} & &=x(t)\\
        &=\dom(x).
    \end{alignat*}
    Para la asociatividad de la acción consideramos,
    \begin{alignat*}{2}
        \dom((gh)x)&=\set{(gh)t\st t\in\dom(x)} 
            &\qquad (gh)x((gh)t)&=gh\cdot x(t)\\
        &=\set{g(ht)\st t\in\dom(x)} & &=g\cdot(h\cdot x(t))\\
        &=\set{gs\st s\in\dom(hx)} & &=g\cdot(hx(ht))\\
        &=\dom(g(hx)) & &=g(hx)(g(ht)).
    \end{alignat*}
    Por ultimo, tenemos que mostrar que la acción es compatible con el valor
    booleano de las atómicas en el sentido de la definición~\ref{def:aest}. Sólo
    veremos que \(\vb{gx\subseteq gy}=g\cdot\vb{x=y}\), ya que el resto se hace
    de la misma manera.
    \begin{align*}
        \vb{gx\subseteq gy}
        &=\bigwedge_{\mathclap{gt\in\dom(gx)}}gx(gt)\Rightarrow\vb{gt\in gy}\\
        &=\bigwedge_{\mathclap{t\in\dom(x)}}g\cdot x(t)\Rightarrow\vb{gt\in gy}
        && \text{definición de \(gx\)}\\
        &=\bigwedge_{\mathclap{t\in\dom(x)}}g\cdot x(t)\Rightarrow g\cdot\vb{t\in y}
        && \text{HI}\\
        &=g\cdot\bigwedge_{\mathclap{t\in\dom(x)}}x(t)\Rightarrow\vb{t\in y}
        && \text{\(g\cdot\) es morfismo}\\
        &=g\cdot\vb{x\subseteq y}. \tag*{\qedhere}
    \end{align*}
\end{proof}

Ahora que ya tenemos que una acción \(G\times B\to B\) por automorfismos induce
una acción de la estructura \(B\)-valuada \(\mb\), podemos ver cómo se comporta
esta acción en algunos elementos.

\begin{lemma}[label=lema:acan]
    Si \(g\in G\) y \(x\in V\), entonces \(g\check{x}=\check{x}\).
\end{lemma}
\begin{proof}
    La demostración es por inducción sobre la relación \(\in\).
    \begin{alignat*}{2}
        \dom(g\check{x})&=\set{gt\st t\in \dom(\check{x})} &\qquad
        g\check{x}(g\check{t})&=g\cdot\check{x}(\check{t})\\
        &=\set{g\check{t}\st t\in x} & &=g\cdot 1\\
        &=\set{\check{t}\st t\in x} & &=1\\
        &=\dom(\check{x}) & &=\check{x}(\check{t})\\
        & & &=\check{x}(g\check{t})
        \tag*{\qedhere}
    \end{alignat*}
\end{proof}

Gracias al lema anterior podemos determinar, de manera sencilla, cómo es el
valor booleano de fórmulas en nombres canónicos.

\begin{lemma}[label=lema:finv]
    Si \(\varphi\) es una fórmula, \(x\in V\) y \(g\in G\), entonces
    \begin{equation*}
        g\cdot\vb{\varphi(\check{x})}=\vb{\varphi(\check{x})}.
    \end{equation*}
\end{lemma}
\begin{proof}
    Esto es una consecuencia inmediata de la proposición~\ref{prop:afor} y el
    lema~\ref{lema:acan}.
\end{proof}

Para dar una consecuencia del lema anterior necesitamos un par de conceptos.

\begin{definition}
    Sean \(B\) un álgebra de Boole y \(G\times B\to B\) una acción por
    automorfismos.
    \begin{enumerate}
        \item Un elemento \(b\in B\) es \emph{invariante} si \(g\cdot b=b\) para
        todo \(g\in G\).
        \item \(B\) es \emph{homogénea} si sus únicos elementos invariantes son
        \(0\) y \(1\).
    \end{enumerate}
\end{definition}

\begin{corollary}
    Si \(B\) es homogénea, \(\varphi\) es una fórmula y \(x\in V\), entonces
    \begin{equation*}
        \vb{\varphi(\check{x})}=0 \quad\text{ó}\quad
        \vb{\varphi(\check{x})}=1.
    \end{equation*}
\end{corollary}
\begin{proof}
    Por el lema~\ref{lema:finv} tenemos que \(\vb{\varphi(\check{x})}\) es un
    elemento invariante del álgebra. Luego, como \(B\) es homogénea se sigue que
    \(\vb{\varphi(\check{x})}=0\) ó \(\vb{\varphi(\check{x})}=1\).
\end{proof}

Dada una acción \(\mb\times G\to G\) nos gustaría decir que un elemento
\(x\in\mb\) es simétrico si \(gx=x\) para \enquote{muchos} \(g\in G\). Con la
terminología de grupos, otra forma de decir esto es que el estabilizador de
\(x\) es un subgrupo grande de \(G\). Recordemos que el estabilizador de \(x\)
es el conjunto
\begin{equation*}
    \stab(x)=\set{g\in G\st gx=x}.
\end{equation*}
Como siempre, la definición de grande está dada por un filtro.

\begin{definition}
    Sea \(G\) un grupo. Decimos que \(\Gamma\) es un filtro de subgrupos de
    \(G\) si para cualesquiera dos subgrupos \(H\) y \(K\) de \(G\) se
    satisface:
    \begin{enumerate}
        \item \(G\in\Gamma\).
        \item Si \(H\in\Gamma\) y \(H\subseteq K\), entonces \(K\in\Gamma\).
        \item Si \(H,K\in\Gamma\), entonces \(H\cap K\in\Gamma\).
    \end{enumerate}
\end{definition}

\begin{definition}
    Sea \(G\) un grupo, \(G\times\mb\to\mb\) una acción y \(\Gamma\) un filtro
    de subgrupos de \(G\). Un elemento \(x\in\mb\)
    es \emph{\(\Gamma\)-simétrico} si \(\stab(x)\in\Gamma\).
\end{definition}

Cuando no haya confunción acerca de qué filtro de subgrupos estamos hablando,
diremos simplemente que \(x\) es simétrico.

Por el lema~\ref{lema:acan} se tiene que \(\stab(\check{x})=G\). Así, los
nombres canónicos de elementos de \(V\) son simétricos.

Ahora podemos dar la estructura de los elementos simétricos. La daremos como una
jerarquía definida por recursión.
\begin{align*}
    &V^{(\Gamma)}_0=\emptyset\\
    &V^{(\Gamma)}_{\alpha+1}
    =\set{x\colon A\to B\st A\subseteq V^{(\Gamma)}_\alpha \text{ y } \stab(x)\in\Gamma}\\
    &V^{(\Gamma)}_{\lambda}=\bigcup_{\alpha<\lambda}V^{(\Gamma)}_{\alpha} 
    &&\text{si \(\lambda\) es límite}\\
    &V^{(\Gamma)}=\bigcup_{\alpha\in\mathrm{Ord}}V^{(\Gamma)}_\alpha.
\end{align*}

Como esta estructura es una subestructura de \(\mb\), podemos tomar el valor
booleano de fórmulas con parámetros en \(V^{(\Gamma)}\) como el mismo que en
\(\mb\). Esto es, si \(\varphi\) es una fórmula y \(x\in \mg\), entonces
\begin{equation*}
    \vb{\varphi(x)}^{\Gamma}=\vb{\varphi(x)}^{B}.
\end{equation*}

De ahora en adelante el objetivo es mostrar que \(V^{(\Gamma)}\) es un modelo de
\textsc{zf}  y que el axioma de elección no se satisface en este modelo. Primero
veamos qué se necesita para que la acción \(G\times\mb\to\mb\) se restrinja a
\(\mg\). Para esto, dado \(x\in\mg\) debemos mostrar que \(gx\in\mg\), es decir,
si \(\stab(x)\in\Gamma\) entonces \(\stab(gx)\in\Gamma\). Para esto, necesitamos
saber cómo se relacionan los estabilizadores de \(x\) y \(gx\).

\begin{lemma}[label=lema:stab]
    Si \(g\in G\) y \(x\in\mb\), entonces \(\stab{gx}=g(\stab(x))g^{-1}\).
\end{lemma}
\begin{proof}
    Notemos que si \(h\in\stab(x)\), entonces \((hg)x=h(gx)=gx\). Si
    multiplicamos ambos lados por \(g^{-1}\) se tiene que \(x=(g^{-1}hg)x\).
    Esto significa que \(g^{-1}hg\in\stab(x)\). Por lo tanto, 
    \(h\in g(\stab(x))g^{-1}\).
    
    Sea \(h\in\stab(x)\) y notemos que
    \begin{equation*}
        (ghg^{-1})(gx)=g(hg^{-1}gx)=g(hx)=gx.
    \end{equation*}
    De esta manera \(ghg^{-1}\in\stab(gx)\).
\end{proof}

Con el lema anterior, para restringir la acción a \(\mg\) necesitamos pedir que
el filtro \(\Gamma\) sea cerrado por conjugación.

\begin{definition}
    Un filtro de subgrupos \(\Gamma\) de \(G\) es \emph{normal}
    si para cualquier \(H\in\Gamma\) y \(g\in G\) se tiene que \(gHg^{-1}\in\Gamma\).
\end{definition}

\begin{corollary}
    Si \(\Gamma\) es un filtro de subgrupos normal de \(G\), entonces la acción
    \(G\times\mb\to\mb\) se restringe a una acción \(G\times\mg\to\mg\).
\end{corollary}
\begin{proof}
    Como \(\Gamma\) es normal, el lema~\ref{lema:stab} implica que si
    \(x\in\mg\) entonces \(gx\in\mg\). Además, como el valor booleano de
    fórmulas con parámetros en \(\mg\) es el mismo que en \(\mb\), se sigue que
    la acción \(G\times\mb\to\mb\) se restringe a una acción
    \(G\times\mg\to\mg\).
\end{proof}

\begin{theorem}
    \(\mg\models\text{\textsc{zf}}\).
\end{theorem}
\begin{proof}
    La
    demostración de extensionalidad y reemplazo son iguales a las de \(\mb\).
    Por el lema~\ref{lema:acan} se tiene que \(\check{\omega}\in\mg\), por lo
    que el axioma de infinito se satisface. Buena fundación se sigue del hecho
    que \(\mg\) es una jerarquía acumulativa, o bien, por ser subclase de
    \(\mb\).

    La demostración del resto de los axiomas seguirá el siguiente esquema: dados
    conjuntos \(x,y\in\mg\) (sólo el axioma del par requiere dos conjuntos, el
    resto requiere uno), se considera el conjunto \(z\in\mb\) que satisface el
    axioma en \(\mb\). Luego, se debe mostrar que \(z\in\mg\), para lo cual se
    demostrará que \(\stab(z)\in\Gamma\). Para esto se usará que
    \(\stab(x),\stab(y)\in\Gamma\) y se demostrará que \(\stab(z)\in\Gamma\).

    Para el axioma del par consideramos \(x,y\in\mg\) y el conjunto en \(\mb\)
    que resulta ser el par de \(x\) e \(y\), es decir, el conjunto
    \begin{equation*}
        \dom(z)=\set{x,y}, \quad z(x)=z(y)=1.
    \end{equation*}
    La demostración consiste en ver que \(\stab(z)\in\Gamma\). Para esto,
    veremos que \(\stab(x)\cap\stab(y)\subseteq\stab(z)\). Sea
    \(g\in\stab(x)\cap\stab(y)\) y notemos que
    \begin{alignat*}{2}
        \dom(gz)&=\set{gx,gy} &\qquad gz(gx)&=g\cdot z(x)\\
        &=\set{x,y} & &=g\cdot 1\\
        &=\dom(z) & &=1=z(gx).
    \end{alignat*}
    Por lo tanto, \(\stab(x)\cap\stab(y)\subseteq\stab(z)\).

    Para el axioma de separación tomemos \(x\in\mg\) y \(\varphi\) una fórmula.
    Recordemos que el conjunto en \(\mb\) que satisface separación está dado por
    \begin{equation*}
        \dom(y)=\dom(x), \quad y(t)=x(t)\land\vb{\varphi(t)}.
    \end{equation*}
    Veamos que \(\stab(x)\subseteq\stab(y)\). Sea \(g\in\stab(x)\). Como
    \(\dom(x)=\dom(gx)\), entonces todo \(s\in\dom(x)\) es de la forma \(gt\)
    para algún \(t\in\dom(x)\).
    \begin{alignat*}{2}
        \dom(gy)&=\set{gt\st t\in\dom(y)} &\qquad gy(s)&=gy(gt)\\
        &=\set{gt\st t\in\dom(x)} & &=g\cdot y(t)\\
        &=\dom(gx) & &=g\cdot(x(t)\land\vb{\varphi(t)})\\
        &=\dom(x) & &=gx(gt)\land\vb{\varphi(gt)}\\
        &=\dom(y) & &=x(s)\land\vb{\varphi(s)}=y(s).
    \end{alignat*}
    Así, \(\stab(x)\subseteq\stab(y)\).
    
    Para el axioma de unión tomemos \(x\in\mg\) y el conjunto en \(\mb\) que
    satisface el axioma de unión, es decir, el conjunto \(y\) dado por
    \begin{equation*}
        \dom(y)=\bigcup_{\mathclap{t\in\dom(x)}}\dom(t), \quad y(s)=1.
    \end{equation*}
    Sea \(g\in\stab(x)\), entonces \(\stab(gy)=\dom(y)\) ya que es la misma
    unión, sólo tiene una permutación de sus uniendos. Además,
    \begin{equation*}
        gy(t)=gy(gs)=g\cdot y(s)=g\cdot 1=1=y(t).
    \end{equation*}
    De esta manera, \(\stab(x)\subseteq\stab(y)\).

    Veamos que se satisface el axioma de potencia. Dado \(x\in\mg\) consideremos
    el conjunto \(y\in\mb\) dado por
    \begin{equation*}
        \dom(y)=\set{u\in\mg\st \dom(u)=\dom(x)}, \quad y(u)=\vb{u\subseteq x}.
    \end{equation*}
    Sea \(g\in\stab(x)\). Si tomamos \(u\in\mg\) tal que \(\dom(u)=\dom(x)\),
    entonces
    \begin{equation*}
        \dom(gu)=\set{gt\st t\in\dom(u)}=\set{gt\st t\in\dom(x)}
        =\dom(gx)=\dom(x)=\dom(u).
    \end{equation*}
    Así, \(\dom(gy)=\dom(y)\). Además,
    \begin{equation*}
        gy(gu)=g\cdot y(u)=g\cdot\vb{u\subseteq x}=\vb{gu\subseteq gx}=\vb{gu\subseteq x}=y(gu).
    \end{equation*}
    Con lo cual concluimos que \(\stab(x)\subseteq\stab(y)\).
\end{proof}

Ahora que ya sabemos que para cualquier grupo \(G\) y cualquier filtro de
subgrupos \(\Gamma\) se tiene que \(\mg\models\text{\textsc{zf}}\), el siguiente
paso es mostrar que el axioma de elección no se satisface en \(\mg\) para
ciertos \(G\), \(\Gamma\) y \(B\).

Lo que haremos para que no se satisfaga el axioma de elección es añadir
\(\omega\) reales y mostrar que el conjunto de estos reales es infinito en el
sentido de Cantor, pero finito en el de Dedekind. Para añadir \(\omega\) reales
sabemos que podemos tomar el orden parcial \(P=\fn(\omega\times\omega,2)\).

Luego, hay que hacer la completación booleana \(e\colon P\to B\). Como \(P\) es
separativo, sólo hay que tomar al espacio topológico \(P\) con la topología del
orden, que denotamos con \(X\), y tomar \(B\) como el álgebra de Boole de los
conjuntos abiertos regulares de este espacio. En este caso, la función \(e\) es
la función que a cada \(p\in P\) le asigna \(\abajo p\).

El siguiente paso es definir una acción en \(B\) por automorfismos. Empezaremos
dando una acción en el orden parcial \(P\), luego la extenderemos al espacio
topológico \(X\) y finalmente al álgebra de Boole \(B\). Definimos una función
\(S_{\omega}\times P\to P\) de la siguiente manera, si \(g\in S_{\omega}\) y
\(p\in P\), entonces
\begin{equation}
    \label{eq:accionP}
    \dom(g^*p)=(\id_{\omega}\times g)^{-1}(\dom(p)), \quad 
    (g^*p)(n,m)=p(n,gm).
\end{equation}

\begin{lemma}
    La función \(S_{\omega}\times P\to P\) de arriba es una acción contravariante.
\end{lemma}
\begin{proof}
    El neutro del grupo \(S_{\omega}\) es la identidad \(\id_{\omega}\). Así, es
    fácil ver que \(\id_{\omega}^*p=p\) para toda \(p\in P\). Para la
    asociatividad consideremos \(g,h\in S_{\omega}\) y \(p\in P\). Entonces,
    \begin{alignat*}{2}
        \dom((gh)^*p)&=(\id_{\omega}\times(gh))^{-1}(\dom(p)) 
            &\qquad (gh)^*p(n,m)&=p(n,(gh)m)\\
        &=(\id_{\omega}\times h)^{-1}((\id_{\omega}\times g)^{-1}(\dom(p)))
        & &=p(n,g(hm))\\
        &=(\id_{\omega}\times h)^{-1}(\dom(g^*p))
        & &=p(n,g(hm))\\
        &=(\id_{\omega}\times h)^{-1}(\dom(g^*p))
        & &=g^*p(n,hm)\\
        &=\dom(h^*(g^*p)) & &=h^*(g^*p)(n,m).
    \end{alignat*}
    Por lo tanto, \((gh)^*p=h^*(g^*p)\).
\end{proof}

\begin{remark}
    Notemos que de la definición de la acción en la ecuación~\eqref{eq:accionP}
    se sigue que si \(q\subseteq p\) entonces \(g^*q\subseteq g^*p\). Esto es,
    \(p\leq q\) implica \(g^*p\leq g^*q\), es decir, la acción preserva el
    orden.
\end{remark}

Ahora veamos que esta acción se extiende al espacio topológico \(X\). Como el
conjunto subyacente de \(X\) es \(P\), entonces la acción de \(S_{\omega}\)
sobre \(X\) se define de la misma manera que sobre \(P\).

\begin{lemma}[label=lema:cont]
    Para cada \(g\in S_{\omega}\) la función \(g^*\colon X\to X\) es continua.
\end{lemma}
\begin{proof}
    Basta mostrar que la preimagen de un abierto básico es un abierto. Sea
    \(p\in P\) y veamos que \((g^*)^{-1}(\abajo p)\) es un abierto. Sean 
    \(q\in (g^*)^{-1}(\abajo p)\) y \(r\leq p\). Como la acción preserva el
    orden tenemos \(g^*r\leq g^*q\in\abajo p\). Por lo tanto, \(r\in
    (g^*)^{-1}(\abajo p)\). Así, \((g^*)^{-1}(\abajo p)\) es un abierto.
\end{proof}

\begin{corollary}
    La acción \(S_{\omega}\times X\to X\) es por homeomorfismos.
\end{corollary}
\begin{proof}
    Como \(S_{\omega}\times P\to P\) es una acción entonces \(g^*\colon P\to P\)
    es biyección, con inversa \((g^{-1})^*\). Por el lema~\ref{lema:cont}
    sabemos que tanto \(g^*\) como \((g^{-1})^*\) son continuas. Por lo tanto,
    \(g^*\) es un homeomorfismo.
\end{proof}

Antes de extender esta acción a \(B\) veamos algunas de las propiedades del
espacio \(X\). Si \(A\subseteq P\), denotaremos con \(\inte(A)\) al interior de
\(A\) y con \(\cer(A)\) a la cerradura de \(A\). Elegimos esta notación para
evitar que se acumulen tantos índices.

\begin{lemma}
    Sea \(A\subseteq P\). Se satisfacen las siguientes condiciones:
    \begin{enumerate}
        \item \(p\in\inte(A)\) si y sólo si \(\abajo p\subseteq A\).
        \item \(\cer(A)=\arriba A=\set{p\in P\st \exists a\in A(a\leq p)}\).
    \end{enumerate}
\end{lemma}
\begin{proof}
    El primer punto es claro ya que \(\abajo p\) es el abierto básico más chico
    que tiene a \(p\).

    Para el segundo punto basta notar que
    \begin{equation*}
        p\in\cer(A) \iff \abajo p\cap A\neq\emptyset \iff 
        \exists a\in A(a\leq p) \iff p\in\arriba A.
        \qedhere
    \end{equation*}
\end{proof}

\end{document}